\documentclass[10pt,a4paper]{amsart}
\usepackage{amsmath,amssymb,amsthm,graphicx}
% Retained preamble for the Stacks Project chapter {CH} (shared project preamble).
% The chapter file pinned-original-source.tex only carries \input{{preamble}} on its first line,
% so no preamble text lives inside the chapter itself.  The Stacks Project uses one single shared
% preamble for every chapter: the snapshot stages it three times (sources/stacks-chow/preamble.tex,
% sources/stacks-etale-cohomology/preamble.tex and focused-seeds/AGP/sources/stacks-snapshot-20261011/
% preamble.tex) and those copies are identical in size.  The block below is copied mechanically,
% line for line, from sources/stacks-chow/preamble.tex lines 14-18 and 154-217: the Xy-pic setup and
% the native theorem declarations together with the mathematical macros.  Nothing was added, removed
% or redefined; the wrapper class and the amsmath/amssymb/amsthm/graphicx base packages are supplied
% by the assembler.
% For commutative diagrams we use Xy-pic
% Chapter: stacks-spaces-more-morphisms.tex -- More on Morphisms of Spaces.
\usepackage[all]{xy}

% We use 2cell for 2-commutative diagrams.
\xyoption{2cell}

% Theorem environments.
%
\theoremstyle{plain}
\newtheorem{theorem}[subsection]{Theorem}
\newtheorem{proposition}[subsection]{Proposition}
\newtheorem{lemma}[subsection]{Lemma}

\theoremstyle{definition}
\newtheorem{definition}[subsection]{Definition}
\newtheorem{example}[subsection]{Example}
\newtheorem{exercise}[subsection]{Exercise}
\newtheorem{situation}[subsection]{Situation}

\theoremstyle{remark}
\newtheorem{remark}[subsection]{Remark}
\newtheorem{remarks}[subsection]{Remarks}

\numberwithin{equation}{subsection}

% Macros
%
\def\lim{\mathop{\mathrm{lim}}\nolimits}
\def\colim{\mathop{\mathrm{colim}}\nolimits}
\def\Spec{\mathop{\mathrm{Spec}}}
\def\Hom{\mathop{\mathrm{Hom}}\nolimits}
\def\Ext{\mathop{\mathrm{Ext}}\nolimits}
\def\SheafHom{\mathop{\mathcal{H}\!\mathit{om}}\nolimits}
\def\SheafExt{\mathop{\mathcal{E}\!\mathit{xt}}\nolimits}
\def\Sch{\mathit{Sch}}
\def\Mor{\mathop{\mathrm{Mor}}\nolimits}
\def\Ob{\mathop{\mathrm{Ob}}\nolimits}
\def\Sh{\mathop{\mathit{Sh}}\nolimits}
\def\NL{\mathop{N\!L}\nolimits}
\def\CH{\mathop{\mathrm{CH}}\nolimits}
\def\proetale{{pro\text{-}\acute{e}tale}}
\def\etale{{\acute{e}tale}}
\def\QCoh{\mathit{QCoh}}
\def\Ker{\mathop{\mathrm{Ker}}}
\def\Im{\mathop{\mathrm{Im}}}
\def\Coker{\mathop{\mathrm{Coker}}}
\def\Coim{\mathop{\mathrm{Coim}}}

% Boxtimes
%
\DeclareMathSymbol{\boxtimes}{\mathbin}{AMSa}{"02}

%
% Macros for moduli stacks/spaces
%
\def\QCohstack{\mathcal{QC}\!\mathit{oh}}
\def\Cohstack{\mathcal{C}\!\mathit{oh}}
\def\Spacesstack{\mathcal{S}\!\mathit{paces}}
\def\Quotfunctor{\mathrm{Quot}}
\def\Hilbfunctor{\mathrm{Hilb}}
\def\Curvesstack{\mathcal{C}\!\mathit{urves}}
\def\Polarizedstack{\mathcal{P}\!\mathit{olarized}}
\def\Complexesstack{\mathcal{C}\!\mathit{omplexes}}
% \Pic is the operator that assigns to X its picard group, usage \Pic(X)
% \Picardstack_{X/B} denotes the Picard stack of X over B
% \Picardfunctor_{X/B} denotes the Picard functor of X over B
\def\Pic{\mathop{\mathrm{Pic}}\nolimits}
\def\Picardstack{\mathcal{P}\!\mathit{ic}}
\def\Picardfunctor{\mathrm{Pic}}
\def\Deformationcategory{\mathcal{D}\!\mathit{ef}}
\usepackage{cleveref}
\crefname{originalreference}{original reference}{original references}
\Crefname{originalreference}{Original reference}{Original references}
\makeatletter
\newlabel{spaces-over-fields-lemma-separably-closed-field-connected-components}{{0A12}{source}{}{sourcepointer-1}{}}
\newlabel{spaces-over-fields-lemma-separably-closed-field-connected-components@cref}{{[originalreference][0][]0A12}{[1][0][]source}{}{}{}}
\makeatother
\title{The count of connected components of geometric fibres is well defined and stable under base change}
\author{The Stacks Project Authors}
\date{stacks-project:spaces-more-morphisms.tex:local-snapshot-20261011; GNU FDL 1.2 or later}
\begin{document}
\maketitle
\noindent\textit{Editorial note.} Counted span is the author's own proof environment, source lines 7086-7108, immediately following the statement at 7069-7084.\par\medskip
\noindent\textit{Editorial note.} One cross-chapter pointer (Spaces over Fields) is recorded in assembly.external\_references.\par\medskip
\noindent\textit{Editorial note.} Same-name pair note: this is the more general side; the schemes side is not counted beside it.\par\medskip
\noindent\textit{Allowed standard prerequisites (editorial).}\begin{itemize}
\item Connected components of a fibre after a separable closure. Spaces over Fields, lemma spaces-over-fields-lemma-separably-closed-field-connected-components, cited at proof lines 7104-7106. Gives bijections between the connected components of the fibres over the various algebraically closed fields representing the same point.
\item Equivalence class of a point of an algebraic space. the definition used in the statement (Spec(k) -> Y in the equivalence class of y with k algebraically closed). Fixes the meaning of the function n\_\{X/Y\} so that well-definedness is a statement about that class.
\item Base change of the fibre. the morphism g : Y' -> Y in the statement and the resulting fibre X'\_\{k'\} at proof lines 7099-7103. Provides the comparison squares whose vertical maps induce the component bijections.
\end{itemize}\medskip
\ifdefined\hypertarget\hypertarget{sourcepointer-1}{}\fi
\noindent\textit{Source reference.} Original reference 0A12: focused-seeds/geometry/S9-spaces-nr-geom-conn-components/ref-tags.txt, source lines 4--4. Navigation pointer only: the counted proof cites this statement by \textbackslash{}ref as a standing input; the deduction the counted proof has to make is carried out inside the counted span, and no part of that deduction is delegated to this pointer.\par\medskip
\begin{lemma}
\label{lemma-base-change-fibres-nr-geometrically-connected-components}
Let $S$ be a scheme.
Let $f : X \to Y$ be a morphism of algebraic spaces over $S$. Let
$$
n_{X/Y} : |Y| \to \{0, 1, 2, 3, \ldots, \infty\}
$$
be the function which associates to $y \in Y$ the number of connected
components of $X_k$ where $\Spec(k) \to Y$ is in the equivalence
class of $y$ with $k$ algebraically closed.
This is well defined and if $g : Y' \to Y$ is a morphism then
$$
n_{X'/Y'} = n_{X/Y} \circ g
$$
where $X' \to Y'$ is the base change of $f$.
\end{lemma}

\begin{proof}
Suppose that $y' \in Y'$ has image $y \in Y$. Let $\Spec(k') \to Y'$
be in the equivalence class of $y'$ with $k'$ algebraically closed.
Then we can choose a commutative diagram
$$
\xymatrix{
\Spec(K) \ar[r] \ar[rd] &
\Spec(k') \ar[r] & Y' \ar[d] \\
& \Spec(k) \ar[r] & Y
}
$$
where $K$ is an algebraically closed field.
The result follows as the morphisms of schemes
$$
\xymatrix{
X'_{k'} & (X'_{k'})_K = (X_k)_K \ar[l] \ar[r] & X_k
}
$$
induce bijections between connected components, see
Spaces over Fields, Lemma
\ref{spaces-over-fields-lemma-separably-closed-field-connected-components}.
To use this to prove the function is well defined take $Y' = Y$.
\end{proof}

\end{document}
